\subsection{The Standard Doclet Tags}
Most of the contents of this chapter is quoted from the document \bdq{JavaDoc Documentation Comment Specification for the Standard Doclet (JDK 25)} \autocite{ORACLE_DOC_JAVADOC_TAG}. 

The output of the Javadoc\index{javadoc} tool is determined by the implementation of the interface \lstinline|jdk.javadoc.doclet.Doclet| \autocite{ORACLE_DOC_DOCLET_INTERFACE} that is specified on the command line for \texttt{javadoc} by the options \texttt{-doclet} and \texttt{-docletpath} \autocite{ORACLE_TOOL:javadoc:options}. The \lstinline|Doclet| implementation is also responsible for interpreting the content of the documentation comments.

If no explicit \lstinline|Doclet| class is specified on the command line, the standard doclet (provided through the class \lstinline|jdk.javadoc.doclet.StandardDoclet| \autocite{ORACLE_DOC_STANDARDDOCLET_CLASS}) is used; it understands the Javadoc tags described below, and it provides an API for additional tags.

Alternative implementations for the \lstinline|Doclet| interface may accept the same syntax as the standard doclet, they may handle the tags completely different or they can ignore them completely. However, due to the support by many tools, the syntax supported by the standard doclet has become a \textit{de facto} standard.

%--------------------------------------------------------------------

\paragraph{\lstinline|@author|}\label{sec:TagAuthor}\index{author@"@author} Usage: \lstinline|@author <name-text>|

The tag adds an \bdq{\textsf{Author:}} entry with the specified name to the generated documentation when the \texttt{-author} option is specified on the command line for the Javadoc\index{javadoc}. A documentation comment can contain multiple \nameref{sec:TagAuthor} tags. Consider to use the custom \bdq{\nameref{sec:TagExtAuthor}}\index{extauthor@"@extauthor}  tag instead.

%--------------------------------------------------------------------

\paragraph{\lstinline|@code|}\label{sec:TagCode}\index{code@"@code}  Usage: \lstinline|{@code <text>}|

This is equivalent to \lstinline|<code>{@literal text}</code>|.

It displays text in the code font without interpreting the text as HTML\index{HTML} markup or nested Javadoc tags. This enables you to use regular angle brackets (\bdq{\textless} and \bdq{\textgreater}) instead of the HTML\index{HTML} entities (\texttt{\&lt;} and \texttt{\&gt;}) in documentation comments, such as in parameter types (\texttt{<Object>}), inequalities ($3 < 4$), or arrows (\texttt{-\textgreater}).

Use the \bdq{\{\nameref{sec:TagLiteral} …\}}\index{literal@"@literal} tag if you want the same functionality without the code font. 

%--------------------------------------------------------------------

\paragraph{\lstinline|@deprecated|}\index{deprecated@"@deprecated}  Usage: \lstinline|@deprecated <text>|

This tag is used in conjunction with the \lstinline|@Deprecated| \autocite{ORACLE_DOC_DEPRECATED_ANNOTATION} annotation to indicate that this API should no longer be used (even though it may continue to work).

The first sentence of the text should tell the user when the API was deprecated and what to use as a replacement. Subsequent sentences can also explain why it was deprecated.

A \bdq{\{\nameref{sec:TagLink} …\}}\index{link@"@link} tag that points to the replacement API should be added where feasible.


%--------------------------------------------------------------------

\paragraph{\lstinline|@docRoot|}  Usage: \lstinline|{@docRoot}|

Represents the relative path to the generated document's (destination) root directory from any generated page. This tag is useful when you want to include a file, such as a copyright page or company logo, that you want to reference from all generated pages.

%--------------------------------------------------------------------

\paragraph{\lstinline|@exception|} This is a synonym for the \bdq{\nameref{sec:TagThrows}}\index{throws@"@throws} tag; it should not be used.

%--------------------------------------------------------------------

\paragraph{\lstinline|@hidden|}\index{hidden@"@hidden}  Usage: \lstinline|@hidden|

Hides a program element from the generated API documentation. This tag may be used when it is not otherwise possible to design the API in a way that such items do not appear at all.

%--------------------------------------------------------------------

\paragraph{\lstinline|@index|}\index{index@"@index}  Usage: \lstinline|{@index <word> <description>}| or\\  
\lstinline|    {@index "<phrase>"  <description>}|

Declares that a word or phrase, together with an optional short description, should appear in the index files generated by the standard doclet. The index entry will be linked to the word or phrase that will appear at this point in the generated documentation. The description may be used when the word or phrase to be indexed is not clear by itself, such as for an acronym.

%--------------------------------------------------------------------

\paragraph{\lstinline|@inheritDoc|}\label{sec:TagInheritDoc}  Usage: \lstinline|{@inheritDoc}|

Inherits (copies) the documentation comment from the nearest inheritable class or implementable interface into the current documentation comment at this tag's location. This enables you to write more general comments higher up the inheritance tree and to write around the copied text.

\paragraph{\lstinline|@link|}\label{sec:TagLink}  Usage: \lstinline|{@link <module/package.class#member> <label>}|

Inserts an inline link with a visible text label that points to the documentation for the specified module, package, class, or member name of a referenced class. 

This tag is similar to the \nameref{sec:TagSee} tag. Both tags require the same references and accept the same syntax for \verb|<module/package.class#member>| and the label. The main difference is that the \lstinline|{@link}| tag generates an inline link rather than placing the link in the “See Also” section. The \lstinline|{@link}| tag begins and ends with curly braces to separate it from the rest of the inline text. If you need to use the right curly brace (“\}”) inside the label, then use the HTML entity notation \verb|&#125;|.

\paragraph{\lstinline|@linkplain|}\label{sec:TagLinkplain}  Usage: \lstinline|{@linkplain <module/package.class#member> <label>}|

Behaves the same as the \nameref{sec:TagLink} tag, except the link label is displayed in plain text rather than code font. Useful when the label is plain text.

\paragraph{\lstinline|@literal|}\label{sec:TagLiteral}  Usage: \lstinline|{@literal <text>}| 

Same as the \nameref{sec:TagCode} tag, but the text is shown as plain text and not in the code font.

\paragraph{\lstinline|@param|}\label{sec:TagParam}  Usage: \lstinline|@param <parameter-name> <description>|

Adds a parameter with the specified parameter name followed by the specified description to the “Parameters” section. The parameter name can be the name of a parameter in a method or constructor, or the name of a type parameter of a class, method, or constructor. Use angle brackets (“\verb#<…>#”) around such a parameter name to indicate the use of a type parameter.

\paragraph{\lstinline|@provides|}\label{sec:TagProvides}  Usage: \lstinline|@provides <service-type> <description>|

This tag may only appear in the documentation comment inside a \verb#module-info.java# file. It serves to document an implementation of a service that is provided by the module. The description may be used to specify how to obtain an instance of this service provider, and any important characteristics of the provider itself. 

\paragraph{\lstinline|@return|}\label{sec:TagReturn}  Usage: \lstinline|@return <description>|

Adds a “Returns” section with the description text to the documentation comment of a method.

\paragraph{\lstinline|@see|}\label{sec:TagSee}  Adds a “See Also” heading with a link or text entry that points to a reference. The \lstinline|@see| tag has three variations; see \autocite{ORACLE_DOC_JAVADOC_TAG} for the details.

\paragraph{\lstinline|@serial|}\label{sec:TagSerial} Used in the documentation comment for a default serializable field. See “Documenting Serializable Fields and Data for a Class”\autocite{ORACLE_DOC_OBJECT_SERIALIZATION:DocumentingSerializableFieldsData}. 

\paragraph{\lstinline|@since|}\label{sec:TagSince}  Usage: \lstinline|@since <since-text>|

Adds a “Since” heading with the specified \verb#<since-text># value to the generated documentation. The text has no special internal structure. This tag that this change or feature has existed since the software release specified by the \verb#<since-text># value, for example: \lstinline|@since 1.5|.

Although it would possible to provide a date or something else, it is recommended to always use a version number with the \lstinline|@since| tag.

\paragraph{\lstinline|@summary|}\label{sec:TagSummary}  Usage:  \lstinline|{@summary <text>}|

Identifies  the summary of an API description, as an alternative to the default policy to identify and use the first sentence of the API description. The tag only has significance when used at the beginning of a description. In all cases, the tag is rendered by simply rendering its content.

The \lstinline|{@summary}| tag has to be used always when a documentation comment has more than one sentence; see also chapter \tqvref{sec:StructureAndContents}.

\paragraph{\lstinline|@throws|}\label{sec:TagThrows}  Usage: \lstinline|@throws <class-name> <description>|

The \lstinline|@throws| tag adds a “Throws” subheading to the generated documentation, with the \verb#<class-name># and the description text. The class name is the name of the exception that might be thrown by the method, and the description provides information about the conditions for that exception to be thrown. 

\paragraph{\lstinline|@uses|}\label{sec:TagUses}  Usage: \lstinline|@uses <service-type> <description>|

This tag may only appear in the documentation comment inside a \verb#module-info.java# file. It serves to document that a service may be used by the module. The description may be used to specify the characteristics of the service that may be required, and what the module will do if no provider for the service is available.

\paragraph{\lstinline|@value|}\label{sec:TagValue}  Usage: \lstinline|{@value}| or \lstinline|{@value <module/package.class#field>}|

This tag is used to display the values of constant in the generated documentation. When the \lstinline|{@value}| tag is used without an argument in the documentation comment of a \lstinline|static final| field, it displays the value of that constant:
\begin{lstlisting}
/**
 * The value of this constant is {@value}.
 */
public static final String SCRIPT_START = "<script>"
\end{lstlisting}

When used with the argument \verb|<module.package.class#field>| in any documentation comment, the \lstinline|{@value}>| tag displays the value of the specified constant:
\begin{lstlisting}
/**
 * Evaluates the script starting with {@value #SCRIPT_START}.
 */
public final String evalScript( String script ) { … }
\end{lstlisting}
The argument \verb|<module.package.class#field>| takes a form similar to that of the \nameref{sec:TagLink}, the \nameref{sec:TagLinkplain}, or the \nameref{sec:TagSee} tag argument, except that the member always must be a \lstinline|static final| field.

\paragraph{\lstinline|@version|}\label{sec:TagVersion}  Usage: \lstinline|@version <version-text>|

Adds a “Version” subheading with the specified \verb#<version-text># value to the generated documents when the \verb#-version# option is specified on the command line. This tag is intended to hold the current release number of the software that this code is part of, as opposed to the \nameref{sec:TagSince} tag, which holds the release number where this code was introduced. The \verb#<version-text># value has no special internal structure.

%==============================================================================
% End of file!!
%==============================================================================

